%% kapa: exam=1 round=37 year=2026 period=1 form=A subject=civil_law subject_ko=민법 seq=1 number=1
%% source: Q-Net 감정평가사 기출문제 (cst00309) artlSeq=5258924 file=raw/1008/5258924/2026년 제37회 감정평가사 1차 1교시.pdf
%% conversion: kordoc 4.24.51 -> tools/gen1.py (auto); review=proofread-sonnet; flags=none
\begin{question}{1}
민법의 법원(法源)에 관한 설명으로 옳은 것은? (다툼이 있으면 판례에 따름)
\begin{choices}
  \item 법률은 관습법에 대하여 열후적ㆍ보충적 성격을 가진다.
  \item 민사에 관한 대통령의 헌법상 긴급명령은 민법의 법원이 될 수 있다.
  \item 법원(法院)은 당사자의 주장이 없으면 직권으로 관습법의 존재를 확정할 수 없다.
  \item 법률행위 해석의 표준인 사실인 관습은 법령으로서의 효력을 가지므로 민법의 법원이다.
  \item 관습법이 그 적용 시점에서 전체 법질서에 부합하지 않게 되었더라도 그 법적 규범으로서의 효력이 부정되지 않는다.
\end{choices}
\end{question}
\answer{2}
